% GENERATED FILE. Edit canonical lessons, then run npm run generate:books.
% canonical-source-hash: fnv1a64:155dace137c5d418
% canonical-lessons: FR-C517-sculpture, FR-C517-sculpter, FR-C517-argile, FR-C517-marbre, FR-C517-sculpteur

\chapter{Marbre et argile}
\label{ch:fr-marbre-et-argile}
\hlchaptermodality{517}

\begin{chapteropening}
\textbf{By the end of this chapter:} \emph{I can talk about sculpture.}

Complete the last lesson of chapter 517: I can talk about sculpture.
\end{chapteropening}

\section[la sculpture]{la sculpture — sculpture}

\label{lesson:FR-C517-sculpture}



\begin{quote}
Before the new one: say the French for easel, then the French for to sketch.
\end{quote}

\subsection*{You'll want to know: la sculpture}
\textbf{la sculpture} — \textquotedblleft{}sculpture\textquotedblright{}.

\begin{quote}
Le musée présente des sculptures modernes.
\end{quote}

\begin{quote}
\emph{The museum is showing modern sculptures.}
\end{quote}

The \emph{p} is silent: \emph{scul-ture}.

\begin{cousinweb}
\emph{la sculpture} is a learned borrowing from Latin \emph{sculptūra} (EtymDB), from \emph{sculpere} \textquotedblleft{}to carve in stone\textquotedblright{} (Lewis \& Short).
\end{cousinweb}

\begin{grammarlens}[title={What you've built}]
One more word for telling, explaining and describing at B1.
\end{grammarlens}

\subsection*{Your turn}
\begin{itemize}
\raggedright
  \item \emph{Say it:} \emph{la sculpture}
  \item \emph{Say it:} \emph{la sculpture}, once more
  \item \emph{Say it:} a sentence of your own with \emph{la sculpture}
  \item \emph{Recall:} say the French for easel, then the French for to sketch, then say \emph{la sculpture} again
\end{itemize}

\begin{tcolorbox}[breakable,colback=teal!4,colframe=teal!35!black,title={Before you move on}]
What does la sculpture mean? (\textquotedblleft{}Sculpture\textquotedblright{}.) Say it once more.
\end{tcolorbox}
\section[sculpter]{sculpter — to sculpt, to carve}

\label{lesson:FR-C517-sculpter}



\begin{quote}
Before the new one: say the French for to sketch, then the French for sculpture.
\end{quote}

\subsection*{You'll want to know: sculpter}
\textbf{sculpter} — \textquotedblleft{}to sculpt, to carve\textquotedblright{}.

\begin{quote}
Il sculpte des animaux dans le bois.
\end{quote}

\begin{quote}
\emph{He carves animals out of wood.}
\end{quote}

The \emph{p} is silent.

\begin{cousinweb}
\emph{sculpter} was made in French from \emph{sculpture} and \emph{sculpteur}, which go back to Latin \emph{sculpere} \textquotedblleft{}to carve\textquotedblright{} (Lewis \& Short).
\end{cousinweb}

\begin{grammarlens}[title={What you've built}]
One more word for telling, explaining and describing at B1.
\end{grammarlens}

\subsection*{Your turn}
\begin{itemize}
\raggedright
  \item \emph{Say it:} \emph{sculpter}
  \item \emph{Say it:} \emph{sculpter}, once more
  \item \emph{Say it:} a sentence of your own with \emph{sculpter}
  \item \emph{Recall:} say the French for to sketch, then the French for sculpture, then say \emph{sculpter} again
\end{itemize}

\begin{tcolorbox}[breakable,colback=teal!4,colframe=teal!35!black,title={Before you move on}]
What does sculpter mean? (\textquotedblleft{}To sculpt, to carve\textquotedblright{}.) Say it once more.
\end{tcolorbox}
\section[l'argile]{l'argile — clay}

\label{lesson:FR-C517-argile}



\begin{quote}
Before the new one: say the French for sculpture, then the French for to sculpt.
\end{quote}

\subsection*{You'll want to know: l'argile}
\textbf{l'argile} — \textquotedblleft{}clay\textquotedblright{}.

\begin{quote}
Les enfants font des petits pots en argile.
\end{quote}

\begin{quote}
\emph{The children are making little pots out of clay.}
\end{quote}

Feminine. \emph{La terre glaise} is a close word.

\begin{cousinweb}
\emph{l'argile} comes from Latin \emph{argilla} (EtymDB), a borrowing from Greek \emph{árgillos} \textquotedblleft{}white clay\textquotedblright{}, related to \emph{argós} \textquotedblleft{}white\textquotedblright{}.
\end{cousinweb}

\begin{grammarlens}[title={What you've built}]
One more word for telling, explaining and describing at B1.
\end{grammarlens}

\subsection*{Your turn}
\begin{itemize}
\raggedright
  \item \emph{Say it:} \emph{l'argile}
  \item \emph{Say it:} \emph{l'argile}, once more
  \item \emph{Say it:} a sentence of your own with \emph{l'argile}
  \item \emph{Recall:} say the French for sculpture, then the French for to sculpt, then say \emph{l'argile} again
\end{itemize}

\begin{tcolorbox}[breakable,colback=teal!4,colframe=teal!35!black,title={Before you move on}]
What does l'argile mean? (\textquotedblleft{}Clay\textquotedblright{}.) Say it once more.
\end{tcolorbox}
\section[le marbre]{le marbre — marble}

\label{lesson:FR-C517-marbre}



\begin{quote}
Before the new one: say the French for to sculpt, then the French for clay.
\end{quote}

\subsection*{You'll want to know: le marbre}
\textbf{le marbre} — \textquotedblleft{}marble\textquotedblright{}.

\begin{quote}
Dans le hall, il y a un escalier en marbre.
\end{quote}

\begin{quote}
\emph{In the hall there is a marble staircase.}
\end{quote}

\emph{Rester de marbre} — to remain unmoved.

\begin{cousinweb}
\emph{le marbre} is inherited from Latin \emph{marmor} \textquotedblleft{}marble\textquotedblright{} (EtymDB; Lewis \& Short), from Greek \emph{mármaros}.
\end{cousinweb}

\begin{grammarlens}[title={What you've built}]
One more word for telling, explaining and describing at B1.
\end{grammarlens}

\subsection*{Your turn}
\begin{itemize}
\raggedright
  \item \emph{Say it:} \emph{le marbre}
  \item \emph{Say it:} \emph{le marbre}, once more
  \item \emph{Say it:} a sentence of your own with \emph{le marbre}
  \item \emph{Recall:} say the French for to sculpt, then the French for clay, then say \emph{le marbre} again
\end{itemize}

\begin{tcolorbox}[breakable,colback=teal!4,colframe=teal!35!black,title={Before you move on}]
What does le marbre mean? (\textquotedblleft{}Marble\textquotedblright{}.) Say it once more.
\end{tcolorbox}
\section[le sculpteur]{le sculpteur — sculptor}

\label{lesson:FR-C517-sculpteur}



\begin{quote}
Before the new one: say the French for clay, then the French for marble.
\end{quote}

\subsection*{You'll want to know: le sculpteur}
\textbf{le sculpteur} — \textquotedblleft{}sculptor\textquotedblright{}.

\begin{quote}
Le sculpteur travaille dans un vieil atelier.
\end{quote}

\begin{quote}
\emph{The sculptor works in an old workshop.}
\end{quote}

A woman is \emph{la sculptrice} or \emph{la sculpteure}. The \emph{p} is silent.

\begin{cousinweb}
\emph{le sculpteur} is a learned borrowing from Latin \emph{sculptor} \textquotedblleft{}a carver\textquotedblright{} (EtymDB), from \emph{sculpere} \textquotedblleft{}to carve\textquotedblright{} (Lewis \& Short).
\end{cousinweb}

\begin{grammarlens}[title={What you've built}]
One more word for telling, explaining and describing at B1.
\end{grammarlens}

\subsection*{Your turn}
\begin{itemize}
\raggedright
  \item \emph{Say it:} \emph{le sculpteur}
  \item \emph{Say it:} \emph{le sculpteur}, once more
  \item \emph{Say it:} a sentence of your own with \emph{le sculpteur}
  \item \emph{Recall:} say the French for clay, then the French for marble, then say \emph{le sculpteur} again
\end{itemize}

\begin{tcolorbox}[breakable,colback=teal!4,colframe=teal!35!black,title={Before you move on}]
What does le sculpteur mean? (\textquotedblleft{}Sculptor\textquotedblright{}.) Say it once more.
\end{tcolorbox}
